\documentclass{article}
\usepackage{lmodern}
\usepackage{amsmath}
\usepackage{listings}
\usepackage{fullpage}
\usepackage{parskip}
\usepackage{xcolor}
\lstset{
	basicstyle=\ttfamily,
	columns=fullflexible,
	frame=single,
	breaklines=true,
	postbreak=\mbox{\textcolor{red}{$\hookrightarrow$}\space},
}
\begin{document}
\title{Experiment `p\_6\_9\_arrays\_equal\_array' Results}
\author{\today}
\date{}
\maketitle
\textbf{Experiment outcome: }SUCCESS
\\\textbf{Bad responses: }0
\\\textbf{Responses containing}~\texttt{assume}~\textbf{: }0
\\\textbf{Responses containing}~\texttt{expect}~\textbf{: }0
\\\textbf{Resolution attempts: }2
\\\textbf{Hard fails (resolution): }1
\\\textbf{Soft fails (resolution): }0
\\\textbf{Verification attempts: }1
\section*{Problem Specification}
\textbf{Problem name: }p\_6\_9\_arrays\_equal\_array
\\\textbf{Natural language statement: }null
\\\textbf{Method signature: }p\_6\_9\_arrays\_equal\_array(first: array<int>, second: array<int>) returns (are\_equal: bool)
\subsection*{Ensures}
\begin{itemize}
\item \texttt{are\_equal == (first.Length == second.Length \&\& forall i :: 0 <= i < first.Length ==> first[i] == second[i])}
\end{itemize}
\clearpage
\section*{GenAI interactions}
Below you will find all interactions between the `user' (program) and the `assistant' (OpenAI).
\subsection*{Program $\rightarrow$ GenAI}
\begin{lstlisting}
You are given the following task to perform in Dafny:

Write a method to satisfy the below requirements.

The signature should be:

p_6_9_arrays_equal_array(first: array<int>, second: array<int>) returns (are_equal: bool)

The method should respect the following contract:

ensures are_equal == (first.Length == second.Length && forall i :: 0 <= i < first.Length ==> first[i] == second[i])

Produce and show only the Dafny body of this method, including the curly braces that surround it. Do not show the signature nor contract. You must not use 'assume' nor 'expect' anywhere in your code.
\end{lstlisting}
\subsection*{GenAI $\rightarrow$ Program}
\textbf{System fingerprint: }null
\\\textbf{ID: }chatcmpl-DnORNw0YgDec9jsrUmeSaSsp726Qq
\\\textbf{Model: }gpt-5.5-2026-04-23
\\\textbf{Created at: }1780664393
\\\textbf{Finish reason: }stop
\\\textbf{Usage: }304 tokens in, and 4470 tokens out
\begin{lstlisting}
{
  if first.Length != second.Length {
    are_equal := false;
    return;
  }

  var i := 0;
  while i < first.Length
    invariant 0 <= i <= first.Length
    invariant first.Length == second.Length
    invariant forall j :: 0 <= j < i ==> first[j] == second[j]
  {
    if first[i] != second[i] {
      are_equal := false;
      assert !(forall j :: 0 <= j < first.Length ==> first[j] == second[j]) by {
        if forall j :: 0 <= j < first.Length ==> first[j] == second[j] {
          assert first[i] == second[i];
          assert false;
        }
      }
      return;
    }

    assert first[i] == second[i];
    assert forall j | 0 <= j < i + 1 :: first[j] == second[j] by {
      if j < i {
        assert first[j] == second[j];
      } else {
        assert j == i;
      }
    }
    i := i + 1;
  }

  assert i == first.Length;
  assert forall j | 0 <= j < first.Length :: first[j] == second[j] by {
    assert 0 <= j < i;
    assert first[j] == second[j];
  }
  are_equal := true;
}
\end{lstlisting}
\subsection*{Program $\rightarrow$ GenAI}
\begin{lstlisting}
When using dafny resolve, the below error is emitted and resolve fails:

problem_data/p_6_9_arrays_equal_array/program.dfy(28,9): Error: unresolved identifier: j
   |
28 |       if j < i {
   |          ^

problem_data/p_6_9_arrays_equal_array/program.dfy(29,21): Error: unresolved identifier: j
   |
29 |         assert first[j] == second[j];
   |                      ^

problem_data/p_6_9_arrays_equal_array/program.dfy(29,34): Error: unresolved identifier: j
   |
29 |         assert first[j] == second[j];
   |                                   ^

problem_data/p_6_9_arrays_equal_array/program.dfy(31,15): Error: unresolved identifier: j
   |
31 |         assert j == i;
   |                ^

problem_data/p_6_9_arrays_equal_array/program.dfy(39,16): Error: unresolved identifier: j
   |
39 |     assert 0 <= j < i;
   |                 ^

problem_data/p_6_9_arrays_equal_array/program.dfy(40,17): Error: unresolved identifier: j
   |
40 |     assert first[j] == second[j];
   |                  ^

problem_data/p_6_9_arrays_equal_array/program.dfy(40,30): Error: unresolved identifier: j
   |
40 |     assert first[j] == second[j];
   |                               ^

7 resolution/type errors detected in program.dfy

Correct the error in code by altering only the method body. Produce and show only the Dafny body, including the curly braces that surround it. Do not show the signature nor contract. You must not use 'assume' nor 'expect' anywhere in your code.
\end{lstlisting}
\subsection*{GenAI $\rightarrow$ Program}
\textbf{System fingerprint: }null
\\\textbf{ID: }chatcmpl-DnOSucvFpOdcSPDMraX44F4TqW3i7
\\\textbf{Model: }gpt-5.5-2026-04-23
\\\textbf{Created at: }1780664488
\\\textbf{Finish reason: }stop
\\\textbf{Usage: }1260 tokens in, and 1981 tokens out
\begin{lstlisting}
{
  if first.Length != second.Length {
    are_equal := false;
    assert first.Length != second.Length;
    return;
  }

  var i := 0;
  while i < first.Length
    invariant 0 <= i <= first.Length
    invariant first.Length == second.Length
    invariant forall j :: 0 <= j < i ==> first[j] == second[j]
  {
    if first[i] != second[i] {
      are_equal := false;
      assert 0 <= i < first.Length;
      assert first[i] != second[i];
      assert !(forall j :: 0 <= j < first.Length ==> first[j] == second[j]) by {
        if forall j :: 0 <= j < first.Length ==> first[j] == second[j] {
          assert first[i] == second[i];
          assert false;
        }
      }
      return;
    }

    assert first[i] == second[i];
    forall j | 0 <= j < i + 1
      ensures first[j] == second[j]
    {
      if j < i {
        assert first[j] == second[j];
      } else {
        assert j == i;
        assert first[j] == second[j];
      }
    }
    assert forall j :: 0 <= j < i + 1 ==> first[j] == second[j];
    i := i + 1;
  }

  assert i == first.Length;
  forall j | 0 <= j < first.Length
    ensures first[j] == second[j]
  {
    assert 0 <= j < i;
    assert first[j] == second[j];
  }
  assert forall j :: 0 <= j < first.Length ==> first[j] == second[j];
  are_equal := true;
}
\end{lstlisting}
\section*{Final Program}
\begin{lstlisting}
method p_6_9_arrays_equal_array(first: array<int>, second: array<int>) returns (are_equal: bool)
	ensures are_equal == (first.Length == second.Length && forall i :: 0 <= i < first.Length ==> first[i] == second[i])
{
  if first.Length != second.Length {
    are_equal := false;
    assert first.Length != second.Length;
    return;
  }

  var i := 0;
  while i < first.Length
    invariant 0 <= i <= first.Length
    invariant first.Length == second.Length
    invariant forall j :: 0 <= j < i ==> first[j] == second[j]
  {
    if first[i] != second[i] {
      are_equal := false;
      assert 0 <= i < first.Length;
      assert first[i] != second[i];
      assert !(forall j :: 0 <= j < first.Length ==> first[j] == second[j]) by {
        if forall j :: 0 <= j < first.Length ==> first[j] == second[j] {
          assert first[i] == second[i];
          assert false;
        }
      }
      return;
    }

    assert first[i] == second[i];
    forall j | 0 <= j < i + 1
      ensures first[j] == second[j]
    {
      if j < i {
        assert first[j] == second[j];
      } else {
        assert j == i;
        assert first[j] == second[j];
      }
    }
    assert forall j :: 0 <= j < i + 1 ==> first[j] == second[j];
    i := i + 1;
  }

  assert i == first.Length;
  forall j | 0 <= j < first.Length
    ensures first[j] == second[j]
  {
    assert 0 <= j < i;
    assert first[j] == second[j];
  }
  assert forall j :: 0 <= j < first.Length ==> first[j] == second[j];
  are_equal := true;
}
\end{lstlisting}
\section*{Total Token Usage}
\textbf{Input tokens: }1564
\\\textbf{Output tokens: }6451
\\\textbf{Reasoning tokens: }5624
\\\textbf{Sum of `total tokens': }8015
\section*{Experiment Timings}
\textbf{Overall Experiment} started at 1780664393510, ended at 1780664556607, lasting 163097ms (163.10 seconds)
\\\textbf{Iteration \#1} started at 1780664393511, ended at 1780664488304, lasting 94793ms (94.79 seconds)
\\\textbf{Iteration \#2} started at 1780664488304, ended at 1780664556607, lasting 68303ms (68.30 seconds)
\end{document}
